\section{Management Plan and Qualifications of the PIs}
\label{sec:pis}

\paragraph{Organization.}
The proposed research is a collaboration between UC Santa Barbara, the lead
organization, and Arizona State University, a subawardee. The PIs at UCSB are Giovanni Vigna, the lead PI, and Christopher
Kruegel, and the PIs at ASU are Adam Doup\'e, who coordinates the subaward,
and Tiffany Bao. Thrust~1 is led at UCSB, drawing on the service-dependency
and network-measurement expertise of PIs Vigna and Kruegel. Thrust~2 is
shared, with UCSB leading the replication and fidelity analysis and ASU the
firmware re-hosting and mock synthesis, and Thrust~3 is led at ASU for the
component-level analysis and at UCSB for the multi-step planning. The
upstream engagement with the anchor projects is coordinated by the lead PI.
The team holds a weekly all-hands meeting, a monthly thrust-level review, and
an annual in-person meeting that rotates between the campuses, and the
students at both institutions share a single repository and continuous
integration system, as they already do on existing joint efforts.

\paragraph{Qualifications.}
The four PIs have a long history of joint research and share a hacking team,
Shellphish, which gives the team the combined expertise in infrastructure
security, program analysis, and autonomous vulnerability discovery that this
project requires. All four have contributed to Shellphish's effort in the
DARPA Artificial Intelligence Cyber Challenge, collaborating on an autonomous
cyber reasoning system that discovers and remediates vulnerabilities in
complex open-source software, the most direct precedent for the agentic
analysis proposed here.

PIs Kruegel and Vigna have a 20-year-long history of collaboration, and they
have developed novel approaches to vulnerability analysis and malware
detection, whose results include open-source tools, such as Anubis and angr,
and hundreds of papers in security and privacy. Both PIs have substantial
experience in managing large grants from NSF, DARPA, and other agencies,
including the ACTION AI Institute, and they co-lead a large lab with dozens
of students. Their previous work on network critical-service
detection~\cite{achilles} and on AI-planning-based privilege
escalation~\cite{chainreactor} underpins Thrusts~1 and~3. PI Doup\'e's
research focuses on automated vulnerability analysis, web security,
cybercrime, and cybersecurity education, and PI Bao's examines the broader
lifecycle of software vulnerabilities, including those in cyber-physical
systems, which bears directly on the fragile-device constraints that shape
Thrust~1.
